\chapter{Los aparatos reproductores}\label{ch:g8:reproductive-systems}

El \cref{ch:g8:sexual-reproduction} necesitaba \omterm{def:g8:sexual-reproduction:gametes}{gametos}; el
\cref{ch:g8:puberty} encendió sus fábricas. Este capítulo cartografía
las propias fábricas: los aparatos reproductores \omterm{def:g4:animal-reproduction:sexual}{masculino} y \omterm{def:g4:animal-reproduction:sexual}{femenino}
--- los órganos, los \omterm{def:g8:sexual-reproduction:gametes}{gametos} que produce cada uno y el notable ritmo
mensual del aparato \omterm{def:g4:animal-reproduction:sexual}{femenino}. Anatomía llana, con nombres llanos: estos
son órganos como el \omterm{def:g4:heart-and-blood:heart}{corazón} y los \omterm{def:g7:how-animals-breathe:four}{pulmones}, y este es su capítulo.

\section{El aparato masculino}

\begin{definition}[El aparato reproductor masculino]\label{def:g8:reproductive-systems:male}
Los órganos del aparato \omterm{def:g4:animal-reproduction:sexual}{masculino}:
\begin{itemize}
  \item los dos \emph{testículos}\index{testículo}, sostenidos fuera
        del calor del cuerpo en el \emph{escroto} --- las fábricas de
        \omterm{def:g8:sexual-reproduction:gametes}{gametos}: desde la \omterm{def:g8:puberty:puberty}{pubertad} producen sin parar
        \emph{espermatozoides}\index{espermatozoide}, por cientos de
        millones;
  \item unos conductos que almacenan y transportan los
        espermatozoides, a los que se unen por el camino unas glándulas
        cuyos líquidos los alimentan y los llevan --- y todo junto forma
        el \emph{semen};
  \item el \emph{pene}\index{pene}, por el que sale el semen del cuerpo
        --- el mismo órgano que, en su otra tarea, lleva la \omterm{def:g7:removing-wastes:kidneys}{orina}
        (\cref{def:g7:removing-wastes:kidneys}); un juego de \omterm{def:g7:heart-and-circulation:rooms}{válvulas}
        mantiene las dos tareas estrictamente separadas.
\end{itemize}
\end{definition}

\begin{example}[El espermatozoide, de cerca]\label{ex:g8:reproductive-systems:sperm}
Al \omterm{def:g5:first-look-at-cells:microscope}{microscopio} (\cref{def:g5:first-look-at-cells:microscope}), un
\omterm{def:g8:reproductive-systems:male}{espermatozoide} es una \omterm{def:g5:first-look-at-cells:cell}{célula} preparada para viajar: una cabeza que es
casi todo \omterm{def:g6:cell-unit-of-life:parts}{núcleo} --- la aportación del padre, empaquetada
(\cref{rem:g5:first-look-at-cells:nucleus}) --- y una cola que se
agita. La \omterm{def:g5:first-look-at-cells:cell}{célula} más pequeña del cuerpo, construida para nadar y
provista para unos días. La producción va de la \omterm{def:g8:puberty:puberty}{pubertad} a la \omterm{def:g3:stages-of-life:stages}{vejez}: la
fábrica \omterm{def:g4:animal-reproduction:sexual}{masculina} es continua.
\end{example}

\section{El aparato femenino}

\begin{definition}[El aparato reproductor femenino]\label{def:g8:reproductive-systems:female}
Los órganos del aparato \omterm{def:g4:animal-reproduction:sexual}{femenino}, todos en la parte baja del vientre:
\begin{itemize}
  \item los dos \emph{ovarios}\index{ovario} --- almacenes y fábricas
        de \omterm{def:g8:sexual-reproduction:gametes}{gametos}: a diferencia de los \omterm{def:g8:reproductive-systems:male}{testículos}, guardan desde el
        \omterm{def:g2:born-grow-die:cycle}{nacimiento} su reserva de futuras \emph{\omterm{def:g5:first-look-at-cells:cell}{células} \omterm{def:g4:animal-reproduction:sexual}{femeninas}} y,
        desde la \omterm{def:g8:puberty:puberty}{pubertad}, las van soltando de una en una;
  \item los dos \emph{oviductos}\index{oviducto}, tubos con flecos que
        atrapan la \omterm{def:g5:first-look-at-cells:cell}{célula} soltada --- el sitio habitual del encuentro
        del paso 2 de la \cref{prop:g8:sexual-reproduction:core};
  \item el \emph{\omterm{prop:g5:human-reproduction:plan}{útero}} (la matriz de la
        \cref{prop:g5:human-reproduction:plan}), cuya pared muscular
        gruesa y cuyo revestimiento renovable son el alojamiento del
        \omterm{prop:g8:sexual-reproduction:core}{cigoto};
  \item la \emph{vagina}, que conecta el \omterm{prop:g5:human-reproduction:plan}{útero} con el exterior ---
        donde se recibe el semen, y el canal del parto del
        \cref{ex:g5:human-reproduction:birth}.
\end{itemize}
\end{definition}

\begin{example}[La célula femenina, de cerca]\label{ex:g8:reproductive-systems:egg}
La \omterm{def:g5:first-look-at-cells:cell}{célula} \omterm{def:g4:animal-reproduction:sexual}{femenina} es lo contrario del \omterm{def:g8:reproductive-systems:male}{espermatozoide} en todo menos en
el papel: la \omterm{def:g5:first-look-at-cells:cell}{célula} \emph{más grande} del cuerpo --- una décima de
milímetro, en el límite de lo que ve el \omterm{def:g1:the-five-senses:sense}{ojo} --- redonda, inmóvil y
cargada de provisiones para los primeros días del viaje (el principio
de la yema de la \cref{rem:g2:animals-grow-up:egg}, a escala de
\omterm{def:g5:first-look-at-cells:cell}{célula}). Una fábrica produce millones sin parar; la otra suelta una
sola \omterm{def:g5:first-look-at-cells:cell}{célula} provista al mes: los números frente a los cuidados de la
\cref{prop:g4:animal-reproduction:strategies}, escritos en los propios
\omterm{def:g8:sexual-reproduction:gametes}{gametos}.
\end{example}

\section{El ciclo}

\begin{proposition}[El ciclo menstrual]\label{prop:g8:reproductive-systems:cycle}
Desde la \omterm{def:g8:puberty:puberty}{pubertad} hasta alrededor de los cincuenta años, el aparato
\omterm{def:g4:animal-reproduction:sexual}{femenino} ejecuta un programa que se repite cada \emph{28 días} más o
menos --- el \emph{ciclo menstrual}\index{ciclo menstrual}, ordenado
por la cadena de \omterm{def:g8:puberty:hormones}{hormonas} de la \cref{def:g8:puberty:hormones}:
\begin{enumerate}
  \item el revestimiento del \omterm{prop:g5:human-reproduction:plan}{útero} se reconstruye, grueso y rico en
        \omterm{prop:g4:journey-of-food:absorb}{sangre} --- el alojamiento preparado;
  \item a mitad del ciclo, alrededor del día 14, un \omterm{def:g8:reproductive-systems:female}{ovario} suelta una
        \omterm{def:g5:first-look-at-cells:cell}{célula} \omterm{def:g4:animal-reproduction:sexual}{femenina} --- la \emph{ovulación}\index{ovulación}; el
        \omterm{def:g8:reproductive-systems:female}{oviducto} la atrapa y durante un día más o menos puede
        fecundarse;
  \item si no hay \omterm{def:g5:flower-to-fruit:steps}{fecundación}: el revestimiento no hace falta; se
        deshace y sale por la vagina a lo largo de unos días --- la
        \emph{menstruación}\index{menstruación} --- y el programa vuelve
        a empezar;
  \item si hay \omterm{def:g5:flower-to-fruit:steps}{fecundación}: el \omterm{prop:g8:sexual-reproduction:core}{cigoto} se instala en el revestimiento
        preparado (\cref{def:g5:human-reproduction:pregnancy}), el ciclo
        se suspende y empieza el \omterm{def:g5:human-reproduction:pregnancy}{embarazo}.
\end{enumerate}
\end{proposition}
\admitted

\begin{omfigure}
\begin{tikzpicture}[scale=0.94]
  % cycle wheel
  \draw[very thick, omDef!60] (0,0) circle (2.3);
  % day marks
  \foreach \a/\d in {90/1, 0/8, -90/14, 180/22} {
    \fill[omDef] (\a:2.3) circle (0.07);
  }
  \node[font=\footnotesize, above] at (90:2.45) {día 1};
  \node[font=\footnotesize, right] at (0:2.45) {día 8};
  \node[font=\footnotesize, below] at (-90:2.45) {día 14};
  \node[font=\footnotesize, left] at (180:2.45) {día 22};
  % phases
  \draw[very thick, omThm] (90:2.3) arc[start angle=90, end angle=35,
    radius=2.3];
  \node[font=\footnotesize, omThm, align=center] at (58:3.5)
    {la menstruación:\\ se desprende el revestimiento};
  \node[font=\footnotesize, omMeth, align=center] at (-35:4.0)
    {revestimiento reconstruido,\\ grueso y listo};
  \fill[omProp] (-90:2.3) circle (0.12);
  \node[font=\footnotesize, omProp, align=center] at (-90:3.55)
    {ovulación:\\ se suelta una célula femenina};
  % center
  \node[font=\small, align=center] at (0,0.25)
    {el ciclo:\\ unos 28 días};
  \node[font=\footnotesize, align=center] at (0,-0.75)
    {(la fecundación lo suspende:\\ embarazo)};
  % arrow direction
  \draw[->, thick, omDef] (150:2.3) arc[start angle=150,
    end angle=120, radius=2.3];
\end{tikzpicture}

{\small El \omterm{prop:g8:reproductive-systems:cycle}{ciclo menstrual} como una rueda: la \omterm{prop:g8:reproductive-systems:cycle}{menstruación} lo abre, el
revestimiento se reconstruye, la \omterm{prop:g8:reproductive-systems:cycle}{ovulación} lo corona a mitad de vuelta
--- y la \omterm{def:g5:flower-to-fruit:steps}{fecundación} es lo único que para la rueda.}
\end{omfigure}

\begin{remark}[Vivir con el ciclo]\label{rem:g8:reproductive-systems:living}
Cuestiones prácticas, dichas con llaneza: los ciclos son irregulares
los primeros años --- la maquinaria se está afinando; la duración varía
de una persona a otra y de un mes a otro; la \omterm{prop:g8:reproductive-systems:cycle}{menstruación} puede traer
dolores (el \omterm{prop:g5:human-reproduction:plan}{útero} es un \omterm{def:g3:bones-and-muscles:muscle}{músculo},
\cref{def:g3:bones-and-muscles:muscle}, y trabaja mientras se
desprende el revestimiento) --- el calor, el movimiento y, cuando haga
falta, el consejo de la enfermera del colegio ayudan; y las compresas,
los tampones y las copas son sencillamente el equipo práctico de la
\omterm{prop:g8:reproductive-systems:cycle}{menstruación}, un pasillo normal de cualquier farmacia. Nada de esto es
una enfermedad: el ciclo es el aparato funcionando exactamente como
está construido.
\end{remark}

\begin{method}[Leer los dos aparatos uno al lado del otro]\label{met:g8:reproductive-systems:sides}
Para no perder el mapa, compara por funciones:
\begin{enumerate}
  \item fábrica de \omterm{def:g8:sexual-reproduction:gametes}{gametos}: \omterm{def:g8:reproductive-systems:male}{testículos} --- \omterm{def:g8:reproductive-systems:female}{ovarios};
  \item \omterm{def:g8:sexual-reproduction:gametes}{gameto}: millones nadadores sin parar --- una \omterm{def:g5:first-look-at-cells:cell}{célula} \omterm{def:g4:animal-reproduction:sexual}{femenina}
        provista al mes;
  \item sitio del encuentro: el \omterm{def:g8:reproductive-systems:female}{oviducto}, al que llega el
        \omterm{def:g8:reproductive-systems:male}{espermatozoide} nadando desde la vagina y a través del \omterm{prop:g5:human-reproduction:plan}{útero};
  \item alojamiento: ninguno en el mapa \omterm{def:g4:animal-reproduction:sexual}{masculino} --- el \omterm{prop:g5:human-reproduction:plan}{útero},
        preparado cada mes, en el \omterm{def:g4:animal-reproduction:sexual}{femenino};
  \item calendario: continuo desde la \omterm{def:g8:puberty:puberty}{pubertad} --- cíclico, de la
        \omterm{def:g8:puberty:puberty}{pubertad} a los cincuenta más o menos.
\end{enumerate}
\end{method}

\section{Ejercicios}

\begin{exercise}[$\star$]\label{exo:g8:reproductive-systems:1}
Cartografía el aparato \omterm{def:g4:animal-reproduction:sexual}{masculino}: nombra cada órgano y su función.
\end{exercise}

\begin{exercise}[$\star$]\label{exo:g8:reproductive-systems:2}
Cartografía igualmente el aparato \omterm{def:g4:animal-reproduction:sexual}{femenino}.
\end{exercise}

\begin{exercise}[$\star$]\label{exo:g8:reproductive-systems:3}
Contrasta el \omterm{def:g8:reproductive-systems:male}{espermatozoide} y la \omterm{def:g5:first-look-at-cells:cell}{célula} \omterm{def:g4:animal-reproduction:sexual}{femenina}: tamaño, construcción,
número y provisiones.
\end{exercise}

\begin{exercise}[$\star$]\label{exo:g8:reproductive-systems:4}
¿Qué pasa en la \omterm{prop:g8:reproductive-systems:cycle}{ovulación}, en qué día aproximado y durante cuánto
tiempo se puede fecundar la \omterm{def:g5:first-look-at-cells:cell}{célula} \omterm{def:g4:animal-reproduction:sexual}{femenina}?
\end{exercise}

\begin{exercise}[$\star$]\label{exo:g8:reproductive-systems:5}
¿Qué es la \omterm{prop:g8:reproductive-systems:cycle}{menstruación}, en términos del revestimiento --- y qué dice
su llegada sobre ese ciclo?
\end{exercise}

\begin{exercise}[$\star$]\label{exo:g8:reproductive-systems:6}
¿Qué acontecimiento suspende el ciclo y qué viene después en el \omterm{prop:g5:human-reproduction:plan}{útero}?
\end{exercise}

\begin{exercise}[$\star\star$]\label{exo:g8:reproductive-systems:7}
¿Dónde ocurren los tres pasos de la
\cref{prop:g8:sexual-reproduction:core} en el mapa de este capítulo ---
órgano por órgano?
\end{exercise}

\begin{exercise}[$\star\star$]\label{exo:g8:reproductive-systems:8}
Lee los \omterm{def:g8:sexual-reproduction:gametes}{gametos} como las dos estrategias de la
\cref{prop:g4:animal-reproduction:strategies} en miniatura. ¿Qué lado
lleva cada lógica?
\end{exercise}

\begin{exercise}[$\star\star$]\label{exo:g8:reproductive-systems:9}
¿Por qué suelen ser irregulares los primeros años de ciclos y por qué
es normal que la duración varíe? (Dos datos de la
\cref{rem:g8:reproductive-systems:living}.)
\end{exercise}

\begin{exercise}[$\star\star$]\label{exo:g8:reproductive-systems:10}
Explica los dolores de la \omterm{prop:g8:reproductive-systems:cycle}{menstruación} con la
\cref{def:g3:bones-and-muscles:muscle} --- y nombra dos cosas que
ayudan.
\end{exercise}

\begin{exercise}[$\star\star$]\label{exo:g8:reproductive-systems:11}
Las dos fábricas llevan calendarios contrarios --- continuo frente a
cíclico y con reserva desde el \omterm{def:g2:born-grow-die:cycle}{nacimiento}. Relaciona cada uno con la
construcción y con el papel de su \omterm{def:g8:sexual-reproduction:gametes}{gameto}.
\end{exercise}

\begin{exercise}[$\star\star\star$]\label{exo:g8:reproductive-systems:12}
``El ciclo es el \omterm{prop:g5:human-reproduction:plan}{útero} haciendo, todos los meses, la pregunta que solo
la \omterm{def:g5:flower-to-fruit:steps}{fecundación} puede responder''. Despliega la imagen en los cuatro
acontecimientos numerados de la
\cref{prop:g8:reproductive-systems:cycle}.
\end{exercise}

\section{Problema: Los modelos didácticos del centro de salud}

\begin{problem}[{Problema de fin de semana --- explicar los aparatos
con dos modelos y un calendario}]\label{pb:g8:reproductive-systems:1}
La jornada de puertas abiertas de un centro de salud usa tres apoyos
didácticos: un modelo del aparato \omterm{def:g4:animal-reproduction:sexual}{masculino}, un modelo del \omterm{def:g4:animal-reproduction:sexual}{femenino} y
un calendario grande de un ciclo de 28 días. Tú eres quien lo explica.

\medskip\noindent\textbf{Parte I --- Los dos modelos.}
\begin{enumerate}
  \item Una visitante recorre el modelo \omterm{def:g4:animal-reproduction:sexual}{masculino} desde los \omterm{def:g8:reproductive-systems:male}{testículos}
        hasta la salida. Narra la ruta, nombrando lo que se le añade
        por el camino y el juego de \omterm{def:g7:heart-and-circulation:rooms}{válvulas} del órgano de doble tarea.
  \item En el modelo \omterm{def:g4:animal-reproduction:sexual}{femenino}, señala dónde espera la reserva de
        \omterm{def:g8:sexual-reproduction:gametes}{gametos}, dónde ocurre normalmente la \omterm{def:g5:flower-to-fruit:steps}{fecundación} y dónde se
        aloja un \omterm{def:g5:human-reproduction:pregnancy}{embarazo}.
  \item Un visitante pregunta por qué se diferencian tanto las fábricas
        de \omterm{def:g8:sexual-reproduction:gametes}{gametos} de los modelos --- millones al día frente a uno al
        mes. Responde con la frase final del
        \cref{ex:g8:reproductive-systems:egg}.
  \item Otra pregunta qué hacían los aparatos antes de la \omterm{def:g8:puberty:puberty}{pubertad}.
        Responde con la maquinaria del \cref{ch:g8:puberty} en una línea.
\end{enumerate}

\medskip\noindent\textbf{Parte II --- El calendario.}
\begin{enumerate}[resume]
  \item Recorre el calendario: coloca la \omterm{prop:g8:reproductive-systems:cycle}{menstruación}, la
        reconstrucción, la \omterm{prop:g8:reproductive-systems:cycle}{ovulación} y el día fértil, con sus fechas
        aproximadas.
  \item Un visitante da por hecho que el día 28 es el mismo para todo
        el mundo. Corrige la suposición con la
        \cref{rem:g8:reproductive-systems:living} --- dos variaciones.
  \item Enseña en el calendario qué cambia si la \omterm{def:g5:flower-to-fruit:steps}{fecundación} ocurre
        cerca del día 14 --- qué acontecimientos se cancelan y cuáles
        empiezan (la \cref{def:g5:human-reproduction:pregnancy}
        continúa la historia).
  \item Una \omterm{def:g3:stages-of-life:stages}{adolescente} pregunta, en voz baja, si los dolores
        significan que algo va mal. Da la respuesta del \omterm{def:g3:bones-and-muscles:muscle}{músculo} y las
        dos ayudas --- y la dirección para cualquier cosa que vaya más
        allá.
\end{enumerate}

\medskip\noindent\textbf{Parte III --- El oficio de explicar.}
\begin{enumerate}[resume]
  \item Los visitantes se ríen con los modelos. Tu frase de apertura
        trata los aparatos como los trata el primer párrafo del
        \cref{ch:g8:reproductive-systems}. Escríbela.
  \item Resume los dos aparatos con la comparación de cinco líneas del
        \cref{met:g8:reproductive-systems:sides}, de memoria.
  \item Un visitante pregunta dónde encajan en los apoyos los comienzos
        de los bebés. Encadena los tres apoyos con los tres pasos de la
        \cref{prop:g8:sexual-reproduction:core}.
  \item Cierra la jornada con una frase sobre por qué enseñan siquiera
        esta anatomía los centros de salud --- conocimiento contra
        leyenda, como argumentaba el \cref{ex:g8:puberty:questions}.
\end{enumerate}
\end{problem}
